\documentclass[11pt,letterpaper]{article}

\usepackage[utf8]{inputenc}
\usepackage[T1]{fontenc}
\usepackage{lmodern}
\usepackage{amsmath,amssymb,amsthm}
\usepackage[margin=1in]{geometry}
\usepackage{hyperref}
\usepackage{microtype}

\title{The BCC Watson Reflection Bridge and the Master Quadratic}
\author{Foundational Ternary Dynamics Project}
\date{\today}

\newtheorem{definition}{Definition}
\newtheorem{theorem}{Theorem}
\newtheorem{corollary}{Corollary}

\begin{document}

\maketitle

\begin{abstract}
We isolate an exact arithmetic bridge between the body-centered cubic (BCC)
Watson integral and the lemniscatic reflection ratio
\[
G^*=\frac{\Gamma(1/4)}{\Gamma(3/4)}.
\]
For the normalized BCC Watson integral,
\[
W_{\rm BCC}=\frac{1}{(2\pi)^3}\int_{[-\pi,\pi]^3}
\frac{d^3k}{1-\cos k_x\cos k_y\cos k_z},
\]
Watson's evaluation and Euler's reflection formula give
\[
W_{\rm BCC}=\frac{\Gamma(1/4)^4}{4\pi^3}
           =\frac{G^{*2}}{2\pi}.
\]
After selecting the normalization
$K:=|\operatorname{Aut}_{\mathbb C}(E,0)|^2(2\pi)W_{\rm BCC}=16G^{*2}$
for $E:y^2=x^3-x$, and after imposing the filtered product rule
$P:=G^*K$, the Vieta data are
\[
x_++x_-=16G^{*2},\qquad x_+x_-=16G^{*3},
\]
and hence the master quadratic
\[
x^2-16G^{*2}x+16G^{*3}=0.
\]
Its larger root is $x_+=137.036171458\ldots$, within $1.26$ ppm of the
measured inverse fine-structure constant. The equality with physical
$\alpha^{-1}$ is not assumed as a theorem here; it is recorded as a
phenomenological selection requiring independent physical matching.
\end{abstract}

\section{The BCC Watson Integral}

We consider the body-diagonal nearest-neighbor walk with steps
$(\pm1,\pm1,\pm1)$ on $\mathbb Z^3$, restricted to a parity class. Geometrically
this is the eight-corner shell of the 26-connected Moore neighborhood, or the
Stella Octangula corner shell; after the usual coordinate scaling it is the BCC
nearest-neighbor pathing structure. The corresponding dimensionless Green
denominator is
\[
1-\cos k_x\cos k_y\cos k_z.
\]

\begin{definition}[Normalized BCC Watson integral]
Define
\[
W_{\rm BCC}=
\frac{1}{(2\pi)^3}\int_{-\pi}^{\pi}\int_{-\pi}^{\pi}\int_{-\pi}^{\pi}
\frac{dk_x\,dk_y\,dk_z}{1-\cos k_x\cos k_y\cos k_z}.
\]
By evenness and periodicity, this equals Watson's first triple integral
\[
\frac{1}{\pi^3}\int_0^\pi\int_0^\pi\int_0^\pi
\frac{du\,dv\,dw}{1-\cos u\cos v\cos w}.
\]
\end{definition}

\begin{theorem}[Watson reflection bridge]
Let
\[
G^*=\frac{\Gamma(1/4)}{\Gamma(3/4)}.
\]
Then
\[
W_{\rm BCC}=\frac{\Gamma(1/4)^4}{4\pi^3}
           =\frac{G^{*2}}{2\pi}.
\]
\end{theorem}

\begin{proof}
The evaluation
\[
W_{\rm BCC}=\frac{\Gamma(1/4)^4}{4\pi^3}
\]
is Watson's first triple integral \cite{Watson1939,MathWorldWatson}; see also
the extended Watson-integral literature \cite{GlasserZucker1977,JoyceZucker2005}.
Euler's reflection formula gives
\[
\Gamma(1/4)\Gamma(3/4)=\pi\sqrt{2}.
\]
This is the $z=1/4$ case of
$\Gamma(z)\Gamma(1-z)=\pi/\sin(\pi z)$ \cite{DLMFReflection}.
Therefore
\[
G^*=\frac{\Gamma(1/4)}{\Gamma(3/4)}
   =\frac{\Gamma(1/4)^2}{\pi\sqrt{2}},
\]
and so
\[
\frac{G^{*2}}{2\pi}
=\frac{\Gamma(1/4)^4}{4\pi^3}
=W_{\rm BCC}.
\]
\end{proof}

\noindent Theorem 1 is classical special-function and lattice-integral
mathematics. Its interpretation as an FTD substrate identity is a model
selection, not a consequence of the formula alone.

\section{Capacity Normalization}

The master quadratic uses the total normalized capacity of a selected BCC/CM
sector. The arithmetic normalization is supplied by the CM elliptic curve
\[
E:y^2=x^3-x,
\]
which has $j(E)=1728$ and complex automorphism group fixing the origin of order
$4$ \cite{Silverman2009}. Thus
\[
|\operatorname{Aut}_{\mathbb C}(E,0)|^2=16.
\]

\begin{definition}[FTD capacity normalization]
For the selected BCC/CM sector, define the total capacity
\[
K=16\cdot(2\pi)\cdot W_{\rm BCC}.
\]
\end{definition}

\noindent This is a normalization rule in the FTD construction. It is not
derived from the Watson integral alone. Substituting Theorem 1 gives
\[
K=16\cdot(2\pi)\cdot\frac{G^{*2}}{2\pi}=16G^{*2}.
\]
Numerically,
\[
K=140.060135374\ldots .
\]

\section{The Master Quadratic}

\begin{definition}[Filtered master polynomial]
Given the capacity $K$ above, define the filtered master polynomial as the
monic quadratic whose elementary symmetric polynomials are
\[
e_1=K,\qquad e_2=G^*K.
\]
Equivalently, if $x_+$ and $x_-$ are its roots, then
\[
x_+ + x_- = K,\qquad x_+x_-=G^*K.
\]
\end{definition}

This definition encodes the self-reference rule of the FTD construction: the
product term is one reflection-ratio step deeper than the trace term.

\begin{corollary}[Master quadratic]
Under Definitions 2 and 3, the filtered master polynomial is
\[
x^2-16G^{*2}x+16G^{*3}=0.
\]
\end{corollary}

\begin{proof}
By Vieta's formulas, the monic quadratic with elementary symmetric data
$e_1=K$ and $e_2=G^*K$ is
\[
x^2-Kx+G^*K=0.
\]
Using Definition 2 and Theorem 1,
\[
K=16G^{*2},\qquad G^*K=16G^{*3}.
\]
Substitution gives the stated polynomial.
\end{proof}

\noindent The Watson/reflection identity is theorem-level. The capacity
normalization and filtered product rule are definitions of the FTD
master-quadratic construction. This separation is essential: the paper proves
the algebraic consequences of the definitions; it does not prove that nature
must choose them.

\section{Roots and Phenomenological Remarks}

The exact roots of the master quadratic are
\[
x_{\pm}=
\frac{16G^{*2}\pm\sqrt{(16G^{*2})^2-64G^{*3}}}{2}.
\]
Equivalently,
\[
x_{\pm}=8G^{*2}\pm4G^{*3/2}\sqrt{4G^*-1}.
\]
Evaluating with
$G^*=2.958675119188638\ldots$ gives
\[
x_+=137.036171458155\ldots,\qquad
x_-=3.023963916339\ldots .
\]

The CODATA 2022 low-energy inverse fine-structure constant is
\[
\alpha(0)^{-1}=137.035999177(21) \quad \cite{NISTAlpha}.
\]
Thus $x_+$ differs from $\alpha(0)^{-1}$ by
\[
1.72281\times 10^{-4},
\]
approximately $1.257$ ppm. This is far outside the CODATA experimental
uncertainty; any physical identification therefore requires additional FTD
matching or correction terms. The mathematical result is exact; the physical
identification
\[
x_+\stackrel{?}{=}\alpha(0)^{-1}
\]
is a phenomenological selection requiring a separate matching from FTD lattice
variables to physical electromagnetic normalization.

The smaller root lies near the integer $3$. In the FTD interpretation, this
motivates the conjectural association with a threefold confined sector,
suggestive of the color rank $N_c=3$. As with the electromagnetic
identification, this is not a theorem of the quadratic alone: the rigorous
statement is the exact value of $x_-$ and its proximity to $3$.

\section{Epistemic Status}

The results above separate into four layers:
\begin{enumerate}
    \item The BCC Watson/reflection identity
    $W_{\rm BCC}=G^{*2}/(2\pi)$ is an exact theorem.
    \item The normalization $K=16\cdot2\pi W_{\rm BCC}$ and filtered product
    $G^*K$ are explicit FTD definitions.
    \item Given these definitions, the master quadratic follows by Vieta's
    formulas.
    \item The identifications $x_+\approx\alpha(0)^{-1}$ and
    $x_-\approx N_c$ are phenomenological selections, not established physical
    derivations in this paper.
\end{enumerate}

\section{Conclusion}

We have isolated the exact arithmetic chain
\[
W_{\rm BCC}=\frac{G^{*2}}{2\pi}
\quad\Longrightarrow\quad
K=16\cdot2\pi W_{\rm BCC}=16G^{*2}
\quad\Longrightarrow\quad
x^2-16G^{*2}x+16G^{*3}=0,
\]
under the stated CM/BCC normalization and filtered self-reference rule. The
larger root is
\[
x_+=137.036171458\ldots,
\]
which is strikingly close to the low-energy physical inverse fine-structure
constant, but not within experimental uncertainty. The paper therefore
establishes a rigorous mathematical source for the master quadratic and records
the physical interpretation as a sharply defined open matching problem rather
than as an unconditional derivation of Standard Model parameters.

\begin{thebibliography}{9}

\bibitem{Watson1939}
G. N. Watson,
``Three triple integrals,''
\emph{Quarterly Journal of Mathematics, Oxford Series} \textbf{10} (1939),
266--276.

\bibitem{GlasserZucker1977}
M. L. Glasser and I. J. Zucker,
``Extended Watson integrals for the cubic lattices,''
\emph{Proceedings of the National Academy of Sciences USA} \textbf{74} (1977),
1800--1801.

\bibitem{JoyceZucker2005}
G. S. Joyce and I. J. Zucker,
``On the evaluation of generalized Watson integrals,''
\emph{Proceedings of the American Mathematical Society} \textbf{133} (2005),
71--81.

\bibitem{DLMFReflection}
NIST Digital Library of Mathematical Functions,
``Gamma Function: Functional Relations,''
\url{https://dlmf.nist.gov/5.5.5}.

\bibitem{Silverman2009}
J. H. Silverman,
\emph{The Arithmetic of Elliptic Curves},
2nd ed., Graduate Texts in Mathematics 106, Springer, 2009.

\bibitem{NISTAlpha}
NIST/CODATA,
``CODATA Value: inverse fine-structure constant,''
\url{https://physics.nist.gov/cgi-bin/cuu/Value?eqalphinv=}.

\bibitem{MathWorldWatson}
E. W. Weisstein,
``Watson's Triple Integrals,''
\emph{MathWorld--A Wolfram Web Resource},
\url{https://mathworld.wolfram.com/WatsonsTripleIntegrals.html}.

\end{thebibliography}

\end{document}
